\section{Binomial}
Una distribución binomial es una distribución discreta que modela el número de éxitos en un número
fijo de ensayos independientes, donde cada ensayo tiene únicamente dos resultados posibles (éxito
o fracaso). En este trabajo, aplicaremos este modelo teórico sobre la variable
\Variable{absence\_days}. Tomaremos como ``éxito'' la condición de asistencia perfecta (0 faltas),
con una probabilidad empírica de \(p = 0{,}3243\) (\(32{,}43\%\)) extraída de nuestra muestra.

El análisis de esta sección se divide en dos partes claramente diferenciadas:

\begin{itemize}
  \item \textbf{Variable aleatoria discreta (VAD):} modelamos empíricamente la variable
  \Variable{absence\_days} (días de ausencia) de nuestra muestra: medidas de tendencia central,
  de dispersión y función de probabilidad.
  \item \textbf{Modelo binomial teórico:} sobre esa VAD aplicamos la distribución binomial
  (éxito = asistencia perfecta, $p = 0{,}3243$) para un escenario hipotético de $n = 10$
  estudiantes seleccionados al azar.
\end{itemize}

\subsection{Análisis de la VAD: \texorpdfstring{\Variable{absence\_days}}{absence\_days}}
A continuación se modela empíricamente la variable \Variable{absence\_days} (días de ausencia) a
partir de las medidas de tendencia central y dispersión y de su función de probabilidad.

\subsubsection{Medidas de tendencia central}
\begin{itemize}
	\item \textbf{Media ($\bar{x}$):} 1,22 días. En promedio, los estudiantes faltan poco más de un día.
	\item \textbf{Mediana ($Me$):} 1 día. El 50\% de los alumnos tiene una falta o menos.
	\item \textbf{Moda ($Mo$):} 0 días. El valor más frecuente es la asistencia perfecta (36 casos).
\end{itemize}

\subsubsection{Medidas de dispersión}
\begin{itemize}
	\item \textbf{Desvío estándar ($S$):} 1,10 días. Los datos se alejan, en promedio, 1,1 unidades de la media.
	\item \textbf{Coeficiente de variación ($CV$):} ~90\%. Al ser mayor al 20\%, indica que los datos son heterogéneos; es decir, hay mucha diferencia entre los que nunca faltan y los que faltan el máximo de días.
\end{itemize}

\subsubsection{Tabla de probabilidades}
Podemos transformar estos datos empíricos en un modelo de probabilidad teórico para predecir el
comportamiento de la población estudiantil:
\begin{DatosTabla}{C{0.14\linewidth} C{0.16\linewidth} C{0.16\linewidth} C{0.16\linewidth} C{0.18\linewidth}}{\textbf{Faltas} & \textbf{P(x)} & \textbf{fi} & \textbf{fr} & \textbf{Fr}}

0 & 32,43\% & 36 & 32,43\% & 32,43\% \\
1 & 27,93\% & 31 & 27,93\% & 60,36\% \\
2 & 24,32\% & 27 & 24,32\% & 84,68\% \\
3 & 15,32\% & 17 & 15,32\% & 100,00\% \\
\textbf{Total} & \textbf{1} & \textbf{111} & \textbf{100,00\%} & \textbf{100,00\%} \\

\end{DatosTabla}


La función de probabilidad $P(x)$ de la variable aleatoria discreta \Variable{absence\_days} es:

% =====================================================================
%  figures/vad-funcion-prob.tex
%  Funcion de probabilidad de la VAD 'absence_days'.
%  (Grafico NUEVO que pide la profesora para el analisis de VAD,
%   distinto del grafico de la distribucion binomial.)
% =====================================================================
\begin{figure}[H]
\centering
\begin{tikzpicture}
\begin{axis}[
  ybar, bar width=22pt,
  title={Función de probabilidad de \Variable{absence\_days} (VAD)},
  xlabel={Días de ausencia ($x$)},
  ylabel={Probabilidad $P(x)$ (\%)},
  ymin=0, ymax=38,
  symbolic x coords={0, 1, 2, 3},
  xtick=data,
  enlarge x limits=0.4,
  nodes near coords,
]
\addplot+[fill=verde, draw=verde] coordinates {(0,32.43) (1,27.93) (2,24.32) (3,15.32)};
\end{axis}
\end{tikzpicture}
\caption{Función de probabilidad de la variable aleatoria discreta \Variable{absence\_days}.}
\label{fig:vad_funcion_probabilidad}
\end{figure}


\newpage
\subsection{Modelo teórico: distribución binomial}
Nos proponemos analizar las probabilidades para un escenario hipotético donde se eligen 10
estudiantes al azar (n = 10), con el objetivo de predecir cuántos de ellos tienen asistencia
perfecta (éxito con probabilidad p = 0,324 o 32,43\%).

\subsubsection{Tabla de análisis binomial}
\begin{DatosTabla}{C{0.05\linewidth} C{0.16\linewidth} L{0.33\linewidth} C{0.14\linewidth} C{0.15\linewidth}}{\textbf{N} & \textbf{Combinaciones} & \textbf{Desarrollo Teórico} & \textbf{Prob. Exacta} & \textbf{Prob. Acumulada}}

0 & 1 & $1 \cdot (0{,}324)^0 \cdot (0{,}676)^{10}$ & 1,98\% & 1,98\% \\
1 & 10 & $10 \cdot (0{,}324)^1 \cdot (0{,}676)^9$ & 9,49\% & 11,47\% \\
2 & 45 & $45 \cdot (0{,}324)^2 \cdot (0{,}676)^8$ & 20,67\% & 32,14\% \\
3 & 120 & $120 \cdot (0{,}324)^3 \cdot (0{,}676)^7$ & 26,33\% & 58,47\% \\
4 & 210 & $210 \cdot (0{,}324)^4 \cdot (0{,}676)^6$ & 22,07\% & 80,54\% \\
5 & 252 & $252 \cdot (0{,}324)^5 \cdot (0{,}676)^5$ & 12,70\% & 93,24\% \\
6 & 210 & $210 \cdot (0{,}324)^6 \cdot (0{,}676)^4$ & 5,06\% & 98,30\% \\
7 & 120 & $120 \cdot (0{,}324)^7 \cdot (0{,}676)^3$ & 1,39\% & 99,69\% \\
8 & 45 & $45 \cdot (0{,}324)^8 \cdot (0{,}676)^2$ & 0,25\% & 99,94\% \\
9 & 10 & $10 \cdot (0{,}324)^9 \cdot (0{,}676)^1$ & 0,03\% & 99,97\% \\
10 & 1 & $1 \cdot (0{,}324)^{10} \cdot (0{,}676)^0$ & 0,00\% & 100\% \\
\textbf{Total} & \textbf{1024} & Suma & \textbf{100\%} & --- \\

\end{DatosTabla}


\begin{itemize}
	\item \textbf{N:} Es la cantidad exacta de alumnos con asistencia perfecta que nos
  interesa encontrar dentro del grupo de 10 (va desde 0 hasta un máximo de 10).
	\item \textbf{Combinaciones:} Indica cuántas formas distintas existen para elegir a esos 'k'
  alumnos dentro del total de 10. Por ejemplo, hay 120 maneras diferentes de combinar a 3 alumnos
  con asistencia perfecta dentro de un grupo de 10.
	\item \textbf{Desarrollo Teórico:} Muestra la aplicación de la fórmula binomial. Multiplica las
  combinaciones por la probabilidad de tener asistencia perfecta elevada a la 'k' $(0.324)^k$, y por
  la probabilidad de NO tener asistencia perfecta elevada al resto del grupo $(0.676)^{10-k}$.
	\item \textbf{Prob. Exacta:} Es el porcentaje de chance que existe de encontrar exactamente
  esa cantidad 'k' de estudiantes. Por ejemplo, hay un 26,33\% de probabilidad de que en el grupo de
  10 alumnos haya exactamente 3 con asistencia perfecta (siendo este el resultado más probable).
	\item \textbf{Prob. Acumulada:} Va sumando las probabilidades exactas fila por fila hacia
  abajo. Sirve para responder preguntas de tipo 'menor o igual a'. Por ejemplo, el valor 0,5847
  (58,47\%) en la fila 3 significa que existe esa probabilidad de encontrar 3 alumnos o menos con
  asistencia perfecta.
\end{itemize}

\subsubsection{Parámetros de la distribución binomial}
\begin{itemize}
\item \textbf{Esperanza Matemática}(Media Teórica) \\
  \PalabraClave{Fórmula:}
  \[ E(X) = \mu = (x_1 * P(x_1)) + (x_2 * P(x_2)) + ... + (x_n * P(x_n)) \]

  \PalabraClave{Cálculo:}
  \begin{equation*}
    \begin{split}
      \sigma^2 = {} & ((0 - 1{,}2253)^2 \cdot 0{,}3243) + ((1 - 1{,}2253)^2 \cdot 0{,}2793) \\
                    & + ((2 - 1{,}2253)^2 \cdot 0{,}2432) + ((3 - 1{,}2253)^2 \cdot 0{,}1532) =
                    1{,}1295 \text{ días}^2
    \end{split}
  \end{equation*}

  \PalabraClave{Cómo se toman los datos:}
  A cada día de ausencia se le resta la esperanza matemática (mu = 1,2253). Ese resultado se eleva al
  cuadrado (para evitar que los valores negativos se cancelen con los positivos) y se multiplica por
  la probabilidad de ese día. Finalmente, se suman los cuatro resultados obtenidos.

  \PalabraClave{Qué significa:} Mide la dispersión o variabilidad de los valores respecto a su media,
  expresada en unidades al cuadrado. Sirve como paso matemático intermedio para conocer la fluctuación
  de los datos escolares de la población.

\item \textbf{Varianza Poblacional} \\
  \PalabraClave{Fórmula:}
  \begin{equation*}
    \sigma^2 = ((x_1 - \mu)^2 \cdot P(x_1)) + ((x_2 - \mu)^2 \cdot P(x_2)) + ... + ((x_n - \mu)^2 \cdot P(x_n))
  \end{equation*}

  \PalabraClave{Desarrollo del cálculo:}
  \begin{equation*}
    \begin{split}
      \sigma^2 = ((0 - 1{,}2253)^2 \cdot 0{,}3243) + ((1 - 1{,}2253)^2 \cdot 0{,}2793) + \\
      ((2 - 1{,}2253)^2 \cdot 0{,}2432) + ((3 - 1{,}2253)^2 \cdot 0{,}1532) = 1{,}1295 \text{ días}^2
    \end{split}
  \end{equation*}

  \PalabraClave{Cómo se toman los datos:}
  A cada día de ausencia se le resta la esperanza matemática (mu = 1,2253). Ese resultado se
  eleva al cuadrado (para evitar que los valores negativos se cancelen con los positivos) y se
  multiplica por la probabilidad de ese día. Finalmente, se suman los cuatro resultados obtenidos.
  
  \PalabraClave{Qué significa:} Mide la dispersión o variabilidad de los valores respecto a su
  media, expresada en unidades al cuadrado. Sirve como paso matemático intermedio para conocer la
  fluctuación de los datos escolares de la población.

\item \textbf{Desvío Estándar Poblacional} \\
  \PalabraClave{Fórmula:}
  \begin{equation*}
    \sigma = \sqrt{\sigma^2}
  \end{equation*}
  \PalabraClave{Desarrollo del cálculo:}
  \begin{equation*}
    \sigma = \sqrt{1{,}1295} \approx 1{,}0632 \text{ días}
  \end{equation*}

  \PalabraClave{Cómo se toman los datos:}
  Se toma el valor final obtenido en el cálculo de la varianza poblacional (1,1295) y se le aplica
  la operación de raíz cuadrada para devolver la medida a su unidad original.

  \PalabraClave{Qué significa:} Indica qué tan dispersos están los datos con respecto al promedio,
  pero expresado directamente en la unidad de medida original (días). En este caso, señala que las
  inasistencias de los alumnos varían, en promedio, 1,06 días alrededor de la media (1,23 días).

\end{itemize}

\subsubsection{Gráfico de barras}
% =====================================================================
%  figures/binomial-barras.tex
%  Distribución binomial (n=10, p=0,324).
%  (Reemplaza a img/barras.png)
% =====================================================================
\begin{figure}[H]
\centering
\begin{tikzpicture}
\begin{axis}[
  ybar, bar width=20pt,
  title={Distribución binomial ($n=10$, $p=0{,}324$)},
  xlabel={Número de estudiantes con asistencia perfecta ($k$)},
  ylabel={Probabilidad (\%)},
  ymin=0, ymax=30,
  xtick=data,
  enlarge x limits=0.08,
]
\addplot+[fill=azul, draw=azul] coordinates {
  (0,1.98) (1,9.49) (2,20.67) (3,26.33) (4,22.07) (5,12.70)
  (6,5.06) (7,1.39) (8,0.25) (9,0.03) (10,0.00)
};
\end{axis}
\end{tikzpicture}
\caption{Gráfico de distribución binomial.}
\label{fig:grafico_probabilidad}
\end{figure}


\PalabraClave{Análisis del Comportamiento de las Barras:}
\begin{itemize}
  \item \textbf{El punto más alto (La Moda):}
    La barra correspondiente a 3 estudiantes es la más alta, alcanzando una altura del 26,33\%. Esto
    demuestra visualmente que el resultado más probable al armar el grupo es encontrar exactamente a
    3 alumnos con asistencia perfecta, lo cual coincide casi perfectamente con la esperanza
    matemática general.

  \item \textbf{Rango central de normalidad:}
    Las barras de 2, 3 y 4 estudiantes concentran la mayor parte de la altura del gráfico. Al estar
    agrupadas en el centro, indican que en casi el 70\% de los casos, la cantidad de alumnos con
    asistencia perfecta se ubicará en este rango.

  \item \textbf{Asimetría de la distribución:}
    El gráfico tiene un ligero sesgo hacia la derecha. Esto se debe a que la probabilidad de éxito
    (32,4\%) es menor al 50\%. Por lo tanto, las barras descienden más lentamente hacia los números
    altos que hacia los bajos.

  \item \textbf{Los extremos insignificantes:}
    A partir de los 7 estudiantes en adelante, las barras se vuelven prácticamente invisibles. Esto
    demuestra de forma gráfica que la probabilidad de encontrar 8, 9 o 10 alumnos con asistencia
    perfecta en un mismo grupo de 10 es casi del 0\%.
\end{itemize}